% plantilla.tex - Plantilla LuaLaTeX para documentos académicos.
% Para usar con pandoc y lualatex.
%
% Los partials tienen el prefijo "articulo-" y viven junto a este archivo.
% Pandoc los resuelve desde el directorio de la plantilla maestra.
%
% Índice de partials:
%   articulo-01-clase.tex               Clase base (article) y opciones YAML
%   articulo-02-fuentes.tex             Motor de fuentes (fontspec, tres familias)
%   articulo-03-idioma.tex              Babel: español, francés, inglés
%   articulo-04-pagina.tex              Geometría de página (geometry)
%   articulo-05-tipografia.tex          Microtype, interlineado, párrafo, viudas
%   articulo-06-romanos.tex             Numeración romana de páginas preliminares
%   articulo-07-encabezado.tex          Encabezado y pie de página (fancyhdr)
%   articulo-08-indice.tex              Índice de contenido (tocloft)
%   articulo-09-colores-enlaces.tex     Colores, hipervínculos, URL, \pandocbounded
%   articulo-10-listas.tex              Listas y tablas (enumitem, booktabs, longtable)
%   articulo-11-citas.tex               Citas exentas (quoting, \quotingbreak)
%   articulo-12-notas.tex               Notas al pie (footmisc, \@makefntext)
%   articulo-13-secciones.tex           Secciones y títulos (needspace, titleformat)
%   articulo-14-figuras-tablas.tex      Gráficos, captions, tablas, matemáticas
%   articulo-15-codigo.tex              Código fuente (fancyvrb, tokens de resaltado)
%   articulo-16-bibliografia.tex        Bibliografía CSL (citeproc/pandoc)
%   articulo-17-abstract.tex            Entorno abstract (Resumen + palabras clave)

% ── Partials ──────────────────────────────────────────────────────────────────

${ articulo-01-clase() }

${ articulo-02-fuentes() }

${ articulo-03-idioma() }

${ articulo-04-pagina() }

${ articulo-05-tipografia() }

${ articulo-06-romanos() }

${ articulo-07-encabezado() }

${ articulo-08-indice() }

${ articulo-09-colores-enlaces() }

${ articulo-10-listas() }

${ articulo-11-citas() }

${ articulo-12-notas() }

${ articulo-13-secciones() }

${ articulo-14-figuras-tablas() }

${ articulo-15-codigo() }

${ articulo-16-bibliografia() }

${ articulo-17-abstract() }

% ── Cabeceras personalizadas de pandoc ────────────────────────────────────────

$for(header-includes)$
$header-includes$
$endfor$

% ── Fecha (mes y año, sin día) ────────────────────────────────────────────────

\newcommand{\mesanio}{%
  \ifcase\month\or enero\or febrero\or marzo\or abril\or mayo\or junio\or
  julio\or agosto\or septiembre\or octubre\or noviembre\or diciembre\fi
  { }de{ }\the\year}

% ── Portada ───────────────────────────────────────────────────────────────────
% \pagenumbering{roman} desde la portada: el contador empieza en i aquí,
% aunque el folio no sea visible hasta el índice (\thispagestyle{empty}).
% \patchcmd suprime el \newpage interno de \maketitle para que
% \thispagestyle{empty} se aplique sobre la portada y no sobre la página siguiente.

$if(title)$
\title{$title$$if(subtitle)$\\\medskip\large $subtitle$$endif$}
$endif$
$if(author)$
\author{$author$}
$endif$
\date{$if(date)$$date$$else$\mesanio$endif$}

\begin{document}

$if(maketitle)$
\pagenumbering{roman}
\makeatletter
\patchcmd{\@maketitle}{\newpage}{}{}{}
\makeatother
\pdfbookmark[1]{Portada}{portada}
\maketitle
\thispagestyle{empty}
$endif$

% ── Resumen y palabras clave ──────────────────────────────────────────────────
% Formato definido en articulo-17-abstract.tex.
% Variables YAML:
%   abstract:  texto del resumen (activa el bloque)
%   keywords:  lista de palabras clave; se imprimen en la misma página

$if(abstract)$
\clearpage
\thispagestyle{empty}
\begin{abstract}
$abstract$
$if(keywords)$
\par\noindent\textbf{Palabras clave:} $for(keywords)$$keywords$$sep$, $endfor$
$endif$
\end{abstract}
$endif$

% ── Dedicatoria ───────────────────────────────────────────────────────────────
% Variable YAML:
%   dedicatoria: texto en cursiva situado en el tercio superior de la página,
%                en un bloque justificado de media caja anclado a la derecha.
%
% \noindent\hfill empuja el minipage al margen derecho.
% Sin \raggedright: el texto queda justificado dentro del bloque.

$if(dedicatoria)$
\clearpage
\thispagestyle{empty}
\vspace*{0.25\textheight}
\noindent\hfill
\begin{minipage}{0.5\linewidth}
  \setstretch{1}%
  \textit{$dedicatoria$}
\end{minipage}
$endif$

% ── Epígrafe ──────────────────────────────────────────────────────────────────
% Variables YAML:
%   epigrafe-texto: cita en redonda (activa el bloque)
%   epigrafe-autor: nombre del autor en versalitas (opcional)
%   epigrafe-obra:  título de la obra en cursiva (opcional)
% Atribución: AUTOR, Obra  (solo los campos declarados)
%
% El texto va en redonda (sin \textit) para distinguirse de la dedicatoria.
% La atribución se alinea a la derecha del minipage con \raggedleft.

$if(epigrafe-texto)$
\clearpage
\thispagestyle{empty}
\vspace*{0.25\textheight}
\noindent\hfill
\begin{minipage}{0.5\linewidth}
  \setstretch{1}%
  $epigrafe-texto$
  $if(epigrafe-autor)$
  \par\smallskip
  {\raggedleft%
  \textsc{$epigrafe-autor$}%
  $if(epigrafe-obra)$, \textit{$epigrafe-obra$}$endif$\par}
  $else$$if(epigrafe-obra)$
  \par\smallskip
  {\raggedleft\textit{$epigrafe-obra$}\par}
  $endif$$endif$
\end{minipage}
$endif$

% ── Índice ────────────────────────────────────────────────────────────────────

$if(toc)$
\clearpage
\renewcommand{\contentsname}{Índice}
% Primera página donde aparece la numeración visible (drop folio).
\thispagestyle{plain}
\pdfbookmark[1]{Índice}{toc}
\tableofcontents
$endif$

% ── Lista de ilustraciones ────────────────────────────────────────────────────
% Variable YAML:
%   lof: true   Activa la lista de figuras (requiere \caption en cada figura).

$if(lof)$
\clearpage
\thispagestyle{plain}
\listoffigures
$endif$

% ── Lista de tablas ───────────────────────────────────────────────────────────
% Variable YAML:
%   lot: true   Activa la lista de tablas (requiere \caption en cada tabla).

$if(lot)$
\clearpage
\thispagestyle{plain}
\listoftables
$endif$

% ── Cuerpo ────────────────────────────────────────────────────────────────────
% impresion: true  →  cuerpo siempre en recto (impar); la página en blanco
%                     lleva numeración romana invisible y arabic arranca en 1
%                     directamente en recto.
% (sin declarar)   →  cuerpo en página nueva, recto o verso según caiga.

$if(impresion)$
\clearpage
\makeatletter
\ifodd\c@page\else
  \hbox{}\thispagestyle{empty}\clearpage
\fi
\makeatother
\pagenumbering{arabic}
$else$
\clearpage
\pagenumbering{arabic}
$endif$

$body$

$for(include-after)$
$include-after$
$endfor$

\end{document}
